% ============================================================
% CAPÍTULO 9: PRINCIPIOS DE BIOENERGÉTICA Y TERMODINÁMICA
% ============================================================
% CONTENIDO BASADO EN: Unidad de Aprendizaje III, Tema "Principios de Bioenergética y termodinámica"
% PROPÓSITO: Identificar las leyes de la termodinámica, explicar reacciones redox, distinguir compuestos de alta energía y clasificar reacciones bioquímicas según su energía libre.
% ============================================================

\section{Leyes de la termodinámica aplicadas a los sistemas biológicos}

% --- CONSEJO: La bioenergética es el estudio de los flujos de energía en los sistemas biológicos. Aplica las leyes de la termodinámica a la vida. ---

\resumebox{definicion}{Bioenergética}{Es la parte de la bioquímica que estudia los procesos de transformación, transferencia y utilización de la energía en los sistemas biológicos.}

\subsection{Primera ley de la termodinámica (Ley de la conservación de la energía)}

% --- CONSEJO: La energía no se crea ni se destruye, solo se transforma. En los seres vivos, la energía química de los alimentos se convierte en otras formas de energía (mecánica, eléctrica, calor). ---
\resumebox{definicion}{Primera ley de la termodinámica}{La energía total de un sistema y su entorno permanece constante. En un sistema biológico, la energía se transforma de una forma a otra, pero no se pierde.}

\subsection{Segunda ley de la termodinámica}

% --- CONSEJO: En cualquier proceso, la entropía (desorden) del universo tiende a aumentar. Los seres vivos son sistemas abiertos que mantienen un orden interno (baja entropía) a expensas de aumentar la entropía del entorno. ---
\resumebox{definicion}{Segunda ley de la termodinámica}{En todo proceso espontáneo, la entropía (desorden) del universo aumenta. Los seres vivos son islas de baja entropía que mantienen su orden consumiendo energía de su entorno.}

\subsection{Energía libre de Gibbs (G)}

% --- CONSEJO: Es la energía disponible para realizar trabajo a temperatura y presión constantes. Es el criterio termodinámico para determinar si una reacción es espontánea. ---
\resumebox{definicion}{Energía libre de Gibbs (G)}{Es la energía que puede ser utilizada para realizar trabajo en un sistema a temperatura y presión constantes. El cambio en la energía libre ($\Delta G$) determina la espontaneidad de una reacción.}

\section{Clasificación de las reacciones bioquímicas}

% --- CONSEJO: Clasifica las reacciones según el cambio de energía libre ($\Delta G$). Esto es fundamental para entender el acoplamiento de reacciones en el metabolismo. ---
\begin{itemize}
    \item \textbf{Reacciones exergónicas:} Son reacciones que liberan energía. Su $\Delta G$ es \textbf{negativo}. Ocurren de forma espontánea. Ejemplo: La hidrólisis del ATP.
    \item \textbf{Reacciones endergónicas:} Son reacciones que requieren energía. Su $\Delta G$ es \textbf{positivo}. No ocurren de forma espontánea y necesitan estar acopladas a una reacción exergónica. Ejemplo: La síntesis de proteínas.
\end{itemize}

% --- CONSEJO: El acoplamiento de reacciones es clave. Las reacciones endergónicas se impulsan acoplándolas a la hidrólisis del ATP (exergónica). ---

\section{Reacciones redox biológicas}

% --- CONSEJO: Explica el concepto de oxidación (pérdida de electrones) y reducción (ganancia de electrones). Estas reacciones son fundamentales en el metabolismo energético. ---
\resumebox{definicion}{Reacción redox}{Es una reacción química en la que hay una transferencia de electrones entre dos especies. La especie que pierde electrones se oxida, y la que gana electrones se reduce.}

\begin{itemize}
    \item En los sistemas biológicos, las reacciones redox están mediadas por coenzimas como \textbf{NAD+ / NADH} y \textbf{FAD / FADH2}.
    \item Estas coenzimas actúan como transportadores de electrones, aceptándolos (reduciéndose) en reacciones catabólicas y donándolos (oxidándose) en otras reacciones.
\end{itemize}

\section{Compuestos de alta energía celular}

% --- CONSEJO: El ATP es la "moneda energética" de la célula. Explica su estructura y por qué es una molécula de alta energía (inestabilidad de los enlaces fosfoanhídrido, repulsión electrostática, estabilización por resonancia de los productos). ---
\resumebox{definicion}{Compuestos de alta energía}{Son moléculas que, al hidrolizarse, liberan una gran cantidad de energía libre ($\Delta G$ es muy negativo). Su hidrólisis se utiliza para impulsar reacciones endergónicas.}

\subsection{Adnosín trifosfato (ATP)}
% --- CONSEJO: El ATP es el principal transportador de energía química en la célula. ---
\begin{itemize}
    \item Está compuesto por: adenina (base nitrogenada), ribosa (azúcar) y tres grupos fosfato.
    \item La energía se almacena en los enlaces fosfoanhídrido entre los grupos fosfato.
    \item La hidrólisis del ATP a ADP + Pi (fosfato inorgánico) libera $\Delta G = -30.5$ kJ/mol en condiciones estándar.
\end{itemize}

\subsection{Guanosín trifosfato (GTP)}
% --- CONSEJO: El GTP también es una molécula de alta energía, utilizada en procesos como la síntesis de proteínas (traducción) y la transducción de señales. ---
\begin{itemize}
    \item Estructura similar al ATP, pero con guanina como base nitrogenada.
    \item La hidrólisis del GTP también libera energía y es utilizada para impulsar reacciones específicas.
\end{itemize}

\subsection{Otros compuestos de alta energía}
\begin{itemize}
    \item \textbf{Fosfocreatina:} Reserva de energía en tejidos musculares y nerviosos. Regenera ATP rápidamente.
    \item \textbf{Acetil-CoA:} Transporta grupos acetilo (de dos carbonos) y es un intermediario clave en el metabolismo.
    \item \textbf{Fosfoenolpiruvato (PEP):} Es un donador de fosfato de alta energía en la glucólisis.
    \item \textbf{1,3-bisfosfoglicerato:} Otro donador de fosfato de alta energía en la glucólisis.
\end{itemize}

\section{Determinación de las condiciones de las reacciones bioquímicas}

% --- CONSEJO: Explica cómo la célula mantiene las condiciones para que las reacciones sean favorables (acoplamiento al ATP, concentraciones de sustratos y productos, etc.). ---
\begin{itemize}
    \item \textbf{Acoplamiento al ATP:} La célula acopla reacciones endergónicas a la hidrólisis del ATP, haciendo que el $\Delta G$ total del proceso sea negativo.
    \item \textbf{Concentración de sustratos y productos:} Las células mantienen una alta relación [ATP] / [ADP] y [NADH] / [NAD+], lo que favorece las reacciones que consumen ATP y las reacciones catabólicas.
    \item \textbf{Localización celular:} Las rutas metabólicas están compartimentadas en diferentes organelos, lo que permite regular su actividad y mantener condiciones óptimas.
    \item \textbf{Regulación enzimática:} La actividad de las enzimas está finamente regulada por inhibidores, activadores y modificaciones covalentes.
\end{itemize}